%% ECON 282E -- Session 5, Deck A6: Worked Example -- the RBC Model by Perturbation
%% The Block A analogue of S03_A4 and S04_B5: the SAME quarterly RBC model,
%% solved a third way, and graded against the Session 4 global solution.
%% Numbers and coordinates from tools/figures/s05_numbers.json, keys
%% "pert_first_order", "higher_order", "three_methods", "fig_irf",
%% "fig_pert_error".

%% ECON 282E -- Foundations of Macroeconomics (UC Riverside, Fall 2026)
%% Shared Beamer preamble for every deck in slides/S01 ... slides/S10.
%%
%% Usage (a deck lives in slides/S0X/ and is compiled from that folder):
%%   %% ECON 282E -- Foundations of Macroeconomics (UC Riverside, Fall 2026)
%% Shared Beamer preamble for every deck in slides/S01 ... slides/S10.
%%
%% Usage (a deck lives in slides/S0X/ and is compiled from that folder):
%%   %% ECON 282E -- Foundations of Macroeconomics (UC Riverside, Fall 2026)
%% Shared Beamer preamble for every deck in slides/S01 ... slides/S10.
%%
%% Usage (a deck lives in slides/S0X/ and is compiled from that folder):
%%   \input{../preamble/course_preamble.tex}
%%   \decktitle{1}{A1}{Who I Am \& Why This Course}{September 24, 2026}
%%   \begin{document}
%%   \makedecktitle
%%   ...
%%
%% Adapted from Courses/AI-ECON-2026/source/shared_preamble.tex (Metropolis theme,
%% concept/caution/example/takeaway boxes, \sectiondivider). Changes for this course:
%%   * course title block (\decktitle / \makedecktitle) instead of \makebeamertitle
%%   * \zh{...} typesets Chinese (book titles) under pdflatex via CJKutf8
%%   * researchbox for the "Research in the wild" frames
%%   * section pages off; TikZ libraries and the course map loaded here
%% Build with pdflatex only (two passes); tools/build_slides.py does this for you.

\documentclass[10pt,english,aspectratio=169]{beamer}

\usepackage[T1]{fontenc}
\usepackage{textcomp}
\usepackage[utf8]{inputenc}
\usepackage{babel}
\usepackage{CJKutf8}

\usepackage{amsmath, amssymb, amsthm, graphicx}
\usepackage{booktabs, multirow, tabularx}
\usepackage{tikz}
\usetikzlibrary{positioning,arrows.meta,shapes,calc,fit,backgrounds}
\usepackage{pgfplots}
\pgfplotsset{compat=1.16}
\usepackage[most]{tcolorbox}
\tcbuselibrary{skins, breakable}
\usepackage{listings}
\usepackage{xcolor}

\lstset{
  basicstyle=\ttfamily\small,
  keywordstyle=\color{blue!80!black}\bfseries,
  commentstyle=\color{green!50!black}\itshape,
  stringstyle=\color{orange!80!black},
  backgroundcolor=\color{gray!8},
  frame=single,
  rulecolor=\color{gray!40},
  breaklines=true,
  showstringspaces=false,
  numbers=none,
}

\usetheme{metropolis}
\usefonttheme{professionalfonts}
\metroset{sectionpage=none}

\definecolor{LightGray}{RGB}{230,230,230}
\definecolor{RedTitle}{RGB}{0,0,200}
\definecolor{VeryLightGray}{RGB}{245,245,245}
\definecolor{DarkText}{RGB}{0,0,0}
\definecolor{UCRBlue}{RGB}{0,45,106}
\definecolor{UCRGold}{RGB}{241,171,0}

% Stage colours of the course arc (used by the course map and elsewhere)
\colorlet{StageTrunk}{blue!12}
\colorlet{StageBranch}{cyan!14}
\colorlet{StageMachine}{gray!18}
\colorlet{StageWall}{red!14}
\colorlet{StageTools}{orange!20}
\colorlet{StageData}{green!16}
\colorlet{StageAgents}{violet!14}

\hypersetup{bookmarks=false,breaklinks=true,pdfborder={0 0 0},colorlinks=true,
  linkcolor=DarkText,urlcolor=RedTitle}

\setbeamercolor{background canvas}{bg=white}
\setbeamercolor{title}{fg=RedTitle}
\setbeamercolor{frametitle}{bg=VeryLightGray,fg=DarkText}
\setbeamercolor{normal text}{fg=DarkText}
\setbeamercolor{alerted text}{fg=RedTitle}
\setbeamersize{text margin left=5mm, text margin right=5mm}
\addtobeamertemplate{frametitle}{}{\vspace{-0.5em}}

\setbeamertemplate{navigation symbols}{}
\setbeamertemplate{footline}{%
  \hspace{5mm}{\usebeamerfont{page number in head/foot}\color{black!45}%
  ECON 282E \,$\cdot$\, \insertshortdeck}\hfill%
  \usebeamercolor[fg]{page number in head/foot}%
  \usebeamerfont{page number in head/foot}%
  \insertframenumber\,/\,\inserttotalframenumber\kern1em\vskip2pt%
}

% ---------------------------------------------------------------------------
% Chinese text under pdflatex (book titles etc.)
\newcommand{\zh}[1]{\begin{CJK*}{UTF8}{gbsn}#1\end{CJK*}}

% ---------------------------------------------------------------------------
% Deck title block
%   \decktitle{session}{deck id}{title}{date}
\newcommand{\insertshortdeck}{}
\newcommand{\decksession}{}
\newcommand{\deckid}{}
\newcommand{\decktitletext}{}
\newcommand{\deckdate}{}
\newcommand{\decktitle}[4]{%
  \renewcommand{\decksession}{#1}%
  \renewcommand{\deckid}{#2}%
  \renewcommand{\decktitletext}{#3}%
  \renewcommand{\deckdate}{#4}%
  \renewcommand{\insertshortdeck}{Session #1 \,$\cdot$\, Deck #1.#2}%
  \title{#3}%
  \author{Zhigang Feng}%
  \date{#4}%
}
\newcommand{\makedecktitle}{%
  \begin{frame}[plain,noframenumbering]
    \begin{tikzpicture}[remember picture,overlay]
      \fill[UCRBlue] (current page.north west) rectangle ([yshift=-0.55cm]current page.north east);
      \fill[UCRGold] ([yshift=-0.55cm]current page.north west) rectangle ([yshift=-0.62cm]current page.north east);
      \node[anchor=west, text=white, font=\small\bfseries] at ([xshift=5mm,yshift=-0.275cm]current page.north west)
        {ECON 282E \,$\cdot$\, Foundations of Macroeconomics};
      \node[anchor=east, text=white, font=\small] at ([xshift=-5mm,yshift=-0.275cm]current page.north east)
        {UC Riverside \,$\cdot$\, Fall 2026};
    \end{tikzpicture}
    \vspace*{1.1cm}
    {\usebeamerfont{title}\color{black!55}\large Session \decksession\ \,$\cdot$\, Deck \decksession.\deckid\par}
    \vspace{0.25cm}
    {\usebeamerfont{title}\color{RedTitle}\LARGE\bfseries \decktitletext\par}
    \vspace{0.7cm}
    {\large Zhigang Feng\par}
    \vspace{0.15cm}
    {\color{black!65}\deckdate\par}
    \vfill
    {\footnotesize\itshape\color{black!55}Build it on paper $\;\to\;$ solve it on the machine $\;\to\;$ push it to the frontier with AI\par}
  \end{frame}}

% ---------------------------------------------------------------------------
% Section divider (plain frame, not counted)
\newcommand{\sectiondivider}[2]{%
\begin{frame}[plain,noframenumbering]
\begin{center}
\vfill
{\Large\textbf{#1}\par}
\vspace{0.35cm}
{\normalsize #2\par}
\vfill
\end{center}
\end{frame}
}

% ---------------------------------------------------------------------------
% Boxes
\newtcolorbox{conceptbox}[1][]{colback=blue!3, colframe=RedTitle!55, boxrule=0.35mm,
  arc=1mm, left=2mm, right=2mm, top=1mm, bottom=1mm, #1}
\newtcolorbox{cautionbox}[1][]{colback=red!3, colframe=red!50!black, boxrule=0.35mm,
  arc=1mm, left=2mm, right=2mm, top=1mm, bottom=1mm, #1}
\newtcolorbox{examplebox}[1][]{colback=green!3, colframe=green!45!black, boxrule=0.35mm,
  arc=1mm, left=2mm, right=2mm, top=1mm, bottom=1mm, #1}
\newtcolorbox{takeawaybox}[1][]{colback=VeryLightGray, colframe=RedTitle!55, boxrule=0.35mm,
  arc=1mm, left=2mm, right=2mm, top=1mm, bottom=1mm, #1}
\newtcolorbox{researchbox}[1][]{enhanced, colback=violet!4, colframe=violet!60!black,
  boxrule=0.35mm, arc=1mm, left=2mm, right=2mm, top=1mm, bottom=1mm,
  fonttitle=\bfseries\small, title={Research in the wild}, #1}

\newcommand{\compacttables}{\renewcommand{\arraystretch}{1.15}}
\newcommand{\src}[1]{{\tiny\color{black!55}#1}}

% ---------------------------------------------------------------------------
\input{../preamble/course_map.tex}

%%   \decktitle{1}{A1}{Who I Am \& Why This Course}{September 24, 2026}
%%   \begin{document}
%%   \makedecktitle
%%   ...
%%
%% Adapted from Courses/AI-ECON-2026/source/shared_preamble.tex (Metropolis theme,
%% concept/caution/example/takeaway boxes, \sectiondivider). Changes for this course:
%%   * course title block (\decktitle / \makedecktitle) instead of \makebeamertitle
%%   * \zh{...} typesets Chinese (book titles) under pdflatex via CJKutf8
%%   * researchbox for the "Research in the wild" frames
%%   * section pages off; TikZ libraries and the course map loaded here
%% Build with pdflatex only (two passes); tools/build_slides.py does this for you.

\documentclass[10pt,english,aspectratio=169]{beamer}

\usepackage[T1]{fontenc}
\usepackage{textcomp}
\usepackage[utf8]{inputenc}
\usepackage{babel}
\usepackage{CJKutf8}

\usepackage{amsmath, amssymb, amsthm, graphicx}
\usepackage{booktabs, multirow, tabularx}
\usepackage{tikz}
\usetikzlibrary{positioning,arrows.meta,shapes,calc,fit,backgrounds}
\usepackage{pgfplots}
\pgfplotsset{compat=1.16}
\usepackage[most]{tcolorbox}
\tcbuselibrary{skins, breakable}
\usepackage{listings}
\usepackage{xcolor}

\lstset{
  basicstyle=\ttfamily\small,
  keywordstyle=\color{blue!80!black}\bfseries,
  commentstyle=\color{green!50!black}\itshape,
  stringstyle=\color{orange!80!black},
  backgroundcolor=\color{gray!8},
  frame=single,
  rulecolor=\color{gray!40},
  breaklines=true,
  showstringspaces=false,
  numbers=none,
}

\usetheme{metropolis}
\usefonttheme{professionalfonts}
\metroset{sectionpage=none}

\definecolor{LightGray}{RGB}{230,230,230}
\definecolor{RedTitle}{RGB}{0,0,200}
\definecolor{VeryLightGray}{RGB}{245,245,245}
\definecolor{DarkText}{RGB}{0,0,0}
\definecolor{UCRBlue}{RGB}{0,45,106}
\definecolor{UCRGold}{RGB}{241,171,0}

% Stage colours of the course arc (used by the course map and elsewhere)
\colorlet{StageTrunk}{blue!12}
\colorlet{StageBranch}{cyan!14}
\colorlet{StageMachine}{gray!18}
\colorlet{StageWall}{red!14}
\colorlet{StageTools}{orange!20}
\colorlet{StageData}{green!16}
\colorlet{StageAgents}{violet!14}

\hypersetup{bookmarks=false,breaklinks=true,pdfborder={0 0 0},colorlinks=true,
  linkcolor=DarkText,urlcolor=RedTitle}

\setbeamercolor{background canvas}{bg=white}
\setbeamercolor{title}{fg=RedTitle}
\setbeamercolor{frametitle}{bg=VeryLightGray,fg=DarkText}
\setbeamercolor{normal text}{fg=DarkText}
\setbeamercolor{alerted text}{fg=RedTitle}
\setbeamersize{text margin left=5mm, text margin right=5mm}
\addtobeamertemplate{frametitle}{}{\vspace{-0.5em}}

\setbeamertemplate{navigation symbols}{}
\setbeamertemplate{footline}{%
  \hspace{5mm}{\usebeamerfont{page number in head/foot}\color{black!45}%
  ECON 282E \,$\cdot$\, \insertshortdeck}\hfill%
  \usebeamercolor[fg]{page number in head/foot}%
  \usebeamerfont{page number in head/foot}%
  \insertframenumber\,/\,\inserttotalframenumber\kern1em\vskip2pt%
}

% ---------------------------------------------------------------------------
% Chinese text under pdflatex (book titles etc.)
\newcommand{\zh}[1]{\begin{CJK*}{UTF8}{gbsn}#1\end{CJK*}}

% ---------------------------------------------------------------------------
% Deck title block
%   \decktitle{session}{deck id}{title}{date}
\newcommand{\insertshortdeck}{}
\newcommand{\decksession}{}
\newcommand{\deckid}{}
\newcommand{\decktitletext}{}
\newcommand{\deckdate}{}
\newcommand{\decktitle}[4]{%
  \renewcommand{\decksession}{#1}%
  \renewcommand{\deckid}{#2}%
  \renewcommand{\decktitletext}{#3}%
  \renewcommand{\deckdate}{#4}%
  \renewcommand{\insertshortdeck}{Session #1 \,$\cdot$\, Deck #1.#2}%
  \title{#3}%
  \author{Zhigang Feng}%
  \date{#4}%
}
\newcommand{\makedecktitle}{%
  \begin{frame}[plain,noframenumbering]
    \begin{tikzpicture}[remember picture,overlay]
      \fill[UCRBlue] (current page.north west) rectangle ([yshift=-0.55cm]current page.north east);
      \fill[UCRGold] ([yshift=-0.55cm]current page.north west) rectangle ([yshift=-0.62cm]current page.north east);
      \node[anchor=west, text=white, font=\small\bfseries] at ([xshift=5mm,yshift=-0.275cm]current page.north west)
        {ECON 282E \,$\cdot$\, Foundations of Macroeconomics};
      \node[anchor=east, text=white, font=\small] at ([xshift=-5mm,yshift=-0.275cm]current page.north east)
        {UC Riverside \,$\cdot$\, Fall 2026};
    \end{tikzpicture}
    \vspace*{1.1cm}
    {\usebeamerfont{title}\color{black!55}\large Session \decksession\ \,$\cdot$\, Deck \decksession.\deckid\par}
    \vspace{0.25cm}
    {\usebeamerfont{title}\color{RedTitle}\LARGE\bfseries \decktitletext\par}
    \vspace{0.7cm}
    {\large Zhigang Feng\par}
    \vspace{0.15cm}
    {\color{black!65}\deckdate\par}
    \vfill
    {\footnotesize\itshape\color{black!55}Build it on paper $\;\to\;$ solve it on the machine $\;\to\;$ push it to the frontier with AI\par}
  \end{frame}}

% ---------------------------------------------------------------------------
% Section divider (plain frame, not counted)
\newcommand{\sectiondivider}[2]{%
\begin{frame}[plain,noframenumbering]
\begin{center}
\vfill
{\Large\textbf{#1}\par}
\vspace{0.35cm}
{\normalsize #2\par}
\vfill
\end{center}
\end{frame}
}

% ---------------------------------------------------------------------------
% Boxes
\newtcolorbox{conceptbox}[1][]{colback=blue!3, colframe=RedTitle!55, boxrule=0.35mm,
  arc=1mm, left=2mm, right=2mm, top=1mm, bottom=1mm, #1}
\newtcolorbox{cautionbox}[1][]{colback=red!3, colframe=red!50!black, boxrule=0.35mm,
  arc=1mm, left=2mm, right=2mm, top=1mm, bottom=1mm, #1}
\newtcolorbox{examplebox}[1][]{colback=green!3, colframe=green!45!black, boxrule=0.35mm,
  arc=1mm, left=2mm, right=2mm, top=1mm, bottom=1mm, #1}
\newtcolorbox{takeawaybox}[1][]{colback=VeryLightGray, colframe=RedTitle!55, boxrule=0.35mm,
  arc=1mm, left=2mm, right=2mm, top=1mm, bottom=1mm, #1}
\newtcolorbox{researchbox}[1][]{enhanced, colback=violet!4, colframe=violet!60!black,
  boxrule=0.35mm, arc=1mm, left=2mm, right=2mm, top=1mm, bottom=1mm,
  fonttitle=\bfseries\small, title={Research in the wild}, #1}

\newcommand{\compacttables}{\renewcommand{\arraystretch}{1.15}}
\newcommand{\src}[1]{{\tiny\color{black!55}#1}}

% ---------------------------------------------------------------------------
%% Course map: the recurring "you are here" slide for ECON 282E.
%% Loaded by course_preamble.tex. Pattern follows AI-ECON-2026 lec01_map.tex, but the map
%% shows the whole ten-session course rather than one lecture.
%%
%%   \CourseMap{s}{label}   s = 1..10 highlights session s and prints "you are here: label"
%%                          (e.g. \CourseMap{1}{1.A2}); earlier sessions get a check mark.
%%   \CourseMap{0}{}        full overview, nothing highlighted.
%%
%% Note: TikZ option lists are comma-split before TeX conditionals run, so every \ifnum
%% below sits outside [...] option lists.

\newcommand{\CourseMap}[2]{%
\begin{center}
\begin{tikzpicture}[
    sess/.style={rectangle, rounded corners=2pt, draw=black!45, minimum width=1.34cm,
        minimum height=1.12cm, text width=1.26cm, align=center, inner sep=1pt},
    stage/.style={font=\tiny\bfseries, text=black!70},
]
  \foreach \i/\d/\t/\c in {%
      1/Sep 24/Intro \\ Growth/StageTrunk,
      2/Oct 1/HA $\cdot$ OLG \\ Policy/StageBranch,
      3/Oct 8/Num.\ DP \\ Python/StageMachine,
      4/Oct 15/Numerics \\ I/StageMachine,
      5/Oct 22/Numerics \\ II/StageMachine,
      6/Oct 29/The wall \\ \textbf{Midterm}/StageWall,
      7/Nov 5/ML \\ Deep nets/StageTools,
      8/Nov 12/RL \\ HA $+$ DL/StageTools,
      9/Nov 19/Text \\ LLMs/StageData,
      10/Dec 3/Agents \\ \textbf{Projects}/StageAgents} {
    \node[sess, fill=\c] (n\i) at ({1.46*(\i-1)}, 0)
      {{\scriptsize\bfseries S\i}\\[-1pt]{\tiny\color{black!60}\d}\\[1pt]{\tiny \t}};
    \ifnum\i<#1
      \node[font=\scriptsize, text=green!45!black] at ([xshift=-2.5mm,yshift=-2.2mm]n\i.north east) {$\checkmark$};
    \fi
    \ifnum\i=#1
      \draw[RedTitle, very thick, rounded corners=2pt]
        ([xshift=-0.6mm,yshift=-0.6mm]n\i.south west) rectangle ([xshift=0.6mm,yshift=0.6mm]n\i.north east);
      % keep the label inside the picture at both ends of the row
      \ifnum\i<3
        \node[font=\scriptsize\bfseries, text=RedTitle, anchor=south west] at ([yshift=1.2mm]n\i.north west)
          {$\blacktriangledown$\,you are here: #2};
      \else\ifnum\i>8
        \node[font=\scriptsize\bfseries, text=RedTitle, anchor=south east] at ([yshift=1.2mm]n\i.north east)
          {you are here: #2\,$\blacktriangledown$};
      \else
        \node[font=\scriptsize\bfseries, text=RedTitle, anchor=south] at ([yshift=1.2mm]n\i.north)
          {you are here: #2\,$\blacktriangledown$};
      \fi\fi
    \fi
  }
  % stage band under the sessions
  \foreach \a/\b/\lab/\c in {1/1/trunk/StageTrunk, 2/2/branches/StageBranch,
      3/5/the machine $\cdot$ classical methods/StageMachine, 6/6/the wall/StageWall,
      7/8/new tools/StageTools, 9/9/new data/StageData, 10/10/colleagues/StageAgents} {
    \fill[\c] ([yshift=-1.5mm]n\a.south west) rectangle ([yshift=-4.6mm]n\b.south east);
    \path ([yshift=-1.5mm]n\a.south west) -- ([yshift=-4.6mm]n\b.south east)
      node[midway, stage] {\lab};
  }
  \node[font=\footnotesize\itshape, text=black!70] at ([yshift=-9.5mm]$(n1.south)!0.5!(n10.south)$)
    {Build it on paper $\;\to\;$ solve it on the machine $\;\to\;$ push it to the frontier with AI};
\end{tikzpicture}
\end{center}
}


%%   \decktitle{1}{A1}{Who I Am \& Why This Course}{September 24, 2026}
%%   \begin{document}
%%   \makedecktitle
%%   ...
%%
%% Adapted from Courses/AI-ECON-2026/source/shared_preamble.tex (Metropolis theme,
%% concept/caution/example/takeaway boxes, \sectiondivider). Changes for this course:
%%   * course title block (\decktitle / \makedecktitle) instead of \makebeamertitle
%%   * \zh{...} typesets Chinese (book titles) under pdflatex via CJKutf8
%%   * researchbox for the "Research in the wild" frames
%%   * section pages off; TikZ libraries and the course map loaded here
%% Build with pdflatex only (two passes); tools/build_slides.py does this for you.

\documentclass[10pt,english,aspectratio=169]{beamer}

\usepackage[T1]{fontenc}
\usepackage{textcomp}
\usepackage[utf8]{inputenc}
\usepackage{babel}
\usepackage{CJKutf8}

\usepackage{amsmath, amssymb, amsthm, graphicx}
\usepackage{booktabs, multirow, tabularx}
\usepackage{tikz}
\usetikzlibrary{positioning,arrows.meta,shapes,calc,fit,backgrounds}
\usepackage{pgfplots}
\pgfplotsset{compat=1.16}
\usepackage[most]{tcolorbox}
\tcbuselibrary{skins, breakable}
\usepackage{listings}
\usepackage{xcolor}

\lstset{
  basicstyle=\ttfamily\small,
  keywordstyle=\color{blue!80!black}\bfseries,
  commentstyle=\color{green!50!black}\itshape,
  stringstyle=\color{orange!80!black},
  backgroundcolor=\color{gray!8},
  frame=single,
  rulecolor=\color{gray!40},
  breaklines=true,
  showstringspaces=false,
  numbers=none,
}

\usetheme{metropolis}
\usefonttheme{professionalfonts}
\metroset{sectionpage=none}

\definecolor{LightGray}{RGB}{230,230,230}
\definecolor{RedTitle}{RGB}{0,0,200}
\definecolor{VeryLightGray}{RGB}{245,245,245}
\definecolor{DarkText}{RGB}{0,0,0}
\definecolor{UCRBlue}{RGB}{0,45,106}
\definecolor{UCRGold}{RGB}{241,171,0}

% Stage colours of the course arc (used by the course map and elsewhere)
\colorlet{StageTrunk}{blue!12}
\colorlet{StageBranch}{cyan!14}
\colorlet{StageMachine}{gray!18}
\colorlet{StageWall}{red!14}
\colorlet{StageTools}{orange!20}
\colorlet{StageData}{green!16}
\colorlet{StageAgents}{violet!14}

\hypersetup{bookmarks=false,breaklinks=true,pdfborder={0 0 0},colorlinks=true,
  linkcolor=DarkText,urlcolor=RedTitle}

\setbeamercolor{background canvas}{bg=white}
\setbeamercolor{title}{fg=RedTitle}
\setbeamercolor{frametitle}{bg=VeryLightGray,fg=DarkText}
\setbeamercolor{normal text}{fg=DarkText}
\setbeamercolor{alerted text}{fg=RedTitle}
\setbeamersize{text margin left=5mm, text margin right=5mm}
\addtobeamertemplate{frametitle}{}{\vspace{-0.5em}}

\setbeamertemplate{navigation symbols}{}
\setbeamertemplate{footline}{%
  \hspace{5mm}{\usebeamerfont{page number in head/foot}\color{black!45}%
  ECON 282E \,$\cdot$\, \insertshortdeck}\hfill%
  \usebeamercolor[fg]{page number in head/foot}%
  \usebeamerfont{page number in head/foot}%
  \insertframenumber\,/\,\inserttotalframenumber\kern1em\vskip2pt%
}

% ---------------------------------------------------------------------------
% Chinese text under pdflatex (book titles etc.)
\newcommand{\zh}[1]{\begin{CJK*}{UTF8}{gbsn}#1\end{CJK*}}

% ---------------------------------------------------------------------------
% Deck title block
%   \decktitle{session}{deck id}{title}{date}
\newcommand{\insertshortdeck}{}
\newcommand{\decksession}{}
\newcommand{\deckid}{}
\newcommand{\decktitletext}{}
\newcommand{\deckdate}{}
\newcommand{\decktitle}[4]{%
  \renewcommand{\decksession}{#1}%
  \renewcommand{\deckid}{#2}%
  \renewcommand{\decktitletext}{#3}%
  \renewcommand{\deckdate}{#4}%
  \renewcommand{\insertshortdeck}{Session #1 \,$\cdot$\, Deck #1.#2}%
  \title{#3}%
  \author{Zhigang Feng}%
  \date{#4}%
}
\newcommand{\makedecktitle}{%
  \begin{frame}[plain,noframenumbering]
    \begin{tikzpicture}[remember picture,overlay]
      \fill[UCRBlue] (current page.north west) rectangle ([yshift=-0.55cm]current page.north east);
      \fill[UCRGold] ([yshift=-0.55cm]current page.north west) rectangle ([yshift=-0.62cm]current page.north east);
      \node[anchor=west, text=white, font=\small\bfseries] at ([xshift=5mm,yshift=-0.275cm]current page.north west)
        {ECON 282E \,$\cdot$\, Foundations of Macroeconomics};
      \node[anchor=east, text=white, font=\small] at ([xshift=-5mm,yshift=-0.275cm]current page.north east)
        {UC Riverside \,$\cdot$\, Fall 2026};
    \end{tikzpicture}
    \vspace*{1.1cm}
    {\usebeamerfont{title}\color{black!55}\large Session \decksession\ \,$\cdot$\, Deck \decksession.\deckid\par}
    \vspace{0.25cm}
    {\usebeamerfont{title}\color{RedTitle}\LARGE\bfseries \decktitletext\par}
    \vspace{0.7cm}
    {\large Zhigang Feng\par}
    \vspace{0.15cm}
    {\color{black!65}\deckdate\par}
    \vfill
    {\footnotesize\itshape\color{black!55}Build it on paper $\;\to\;$ solve it on the machine $\;\to\;$ push it to the frontier with AI\par}
  \end{frame}}

% ---------------------------------------------------------------------------
% Section divider (plain frame, not counted)
\newcommand{\sectiondivider}[2]{%
\begin{frame}[plain,noframenumbering]
\begin{center}
\vfill
{\Large\textbf{#1}\par}
\vspace{0.35cm}
{\normalsize #2\par}
\vfill
\end{center}
\end{frame}
}

% ---------------------------------------------------------------------------
% Boxes
\newtcolorbox{conceptbox}[1][]{colback=blue!3, colframe=RedTitle!55, boxrule=0.35mm,
  arc=1mm, left=2mm, right=2mm, top=1mm, bottom=1mm, #1}
\newtcolorbox{cautionbox}[1][]{colback=red!3, colframe=red!50!black, boxrule=0.35mm,
  arc=1mm, left=2mm, right=2mm, top=1mm, bottom=1mm, #1}
\newtcolorbox{examplebox}[1][]{colback=green!3, colframe=green!45!black, boxrule=0.35mm,
  arc=1mm, left=2mm, right=2mm, top=1mm, bottom=1mm, #1}
\newtcolorbox{takeawaybox}[1][]{colback=VeryLightGray, colframe=RedTitle!55, boxrule=0.35mm,
  arc=1mm, left=2mm, right=2mm, top=1mm, bottom=1mm, #1}
\newtcolorbox{researchbox}[1][]{enhanced, colback=violet!4, colframe=violet!60!black,
  boxrule=0.35mm, arc=1mm, left=2mm, right=2mm, top=1mm, bottom=1mm,
  fonttitle=\bfseries\small, title={Research in the wild}, #1}

\newcommand{\compacttables}{\renewcommand{\arraystretch}{1.15}}
\newcommand{\src}[1]{{\tiny\color{black!55}#1}}

% ---------------------------------------------------------------------------
%% Course map: the recurring "you are here" slide for ECON 282E.
%% Loaded by course_preamble.tex. Pattern follows AI-ECON-2026 lec01_map.tex, but the map
%% shows the whole ten-session course rather than one lecture.
%%
%%   \CourseMap{s}{label}   s = 1..10 highlights session s and prints "you are here: label"
%%                          (e.g. \CourseMap{1}{1.A2}); earlier sessions get a check mark.
%%   \CourseMap{0}{}        full overview, nothing highlighted.
%%
%% Note: TikZ option lists are comma-split before TeX conditionals run, so every \ifnum
%% below sits outside [...] option lists.

\newcommand{\CourseMap}[2]{%
\begin{center}
\begin{tikzpicture}[
    sess/.style={rectangle, rounded corners=2pt, draw=black!45, minimum width=1.34cm,
        minimum height=1.12cm, text width=1.26cm, align=center, inner sep=1pt},
    stage/.style={font=\tiny\bfseries, text=black!70},
]
  \foreach \i/\d/\t/\c in {%
      1/Sep 24/Intro \\ Growth/StageTrunk,
      2/Oct 1/HA $\cdot$ OLG \\ Policy/StageBranch,
      3/Oct 8/Num.\ DP \\ Python/StageMachine,
      4/Oct 15/Numerics \\ I/StageMachine,
      5/Oct 22/Numerics \\ II/StageMachine,
      6/Oct 29/The wall \\ \textbf{Midterm}/StageWall,
      7/Nov 5/ML \\ Deep nets/StageTools,
      8/Nov 12/RL \\ HA $+$ DL/StageTools,
      9/Nov 19/Text \\ LLMs/StageData,
      10/Dec 3/Agents \\ \textbf{Projects}/StageAgents} {
    \node[sess, fill=\c] (n\i) at ({1.46*(\i-1)}, 0)
      {{\scriptsize\bfseries S\i}\\[-1pt]{\tiny\color{black!60}\d}\\[1pt]{\tiny \t}};
    \ifnum\i<#1
      \node[font=\scriptsize, text=green!45!black] at ([xshift=-2.5mm,yshift=-2.2mm]n\i.north east) {$\checkmark$};
    \fi
    \ifnum\i=#1
      \draw[RedTitle, very thick, rounded corners=2pt]
        ([xshift=-0.6mm,yshift=-0.6mm]n\i.south west) rectangle ([xshift=0.6mm,yshift=0.6mm]n\i.north east);
      % keep the label inside the picture at both ends of the row
      \ifnum\i<3
        \node[font=\scriptsize\bfseries, text=RedTitle, anchor=south west] at ([yshift=1.2mm]n\i.north west)
          {$\blacktriangledown$\,you are here: #2};
      \else\ifnum\i>8
        \node[font=\scriptsize\bfseries, text=RedTitle, anchor=south east] at ([yshift=1.2mm]n\i.north east)
          {you are here: #2\,$\blacktriangledown$};
      \else
        \node[font=\scriptsize\bfseries, text=RedTitle, anchor=south] at ([yshift=1.2mm]n\i.north)
          {you are here: #2\,$\blacktriangledown$};
      \fi\fi
    \fi
  }
  % stage band under the sessions
  \foreach \a/\b/\lab/\c in {1/1/trunk/StageTrunk, 2/2/branches/StageBranch,
      3/5/the machine $\cdot$ classical methods/StageMachine, 6/6/the wall/StageWall,
      7/8/new tools/StageTools, 9/9/new data/StageData, 10/10/colleagues/StageAgents} {
    \fill[\c] ([yshift=-1.5mm]n\a.south west) rectangle ([yshift=-4.6mm]n\b.south east);
    \path ([yshift=-1.5mm]n\a.south west) -- ([yshift=-4.6mm]n\b.south east)
      node[midway, stage] {\lab};
  }
  \node[font=\footnotesize\itshape, text=black!70] at ([yshift=-9.5mm]$(n1.south)!0.5!(n10.south)$)
    {Build it on paper $\;\to\;$ solve it on the machine $\;\to\;$ push it to the frontier with AI};
\end{tikzpicture}
\end{center}
}


\graphicspath{{./figures/}}

\lstset{language=Python, basicstyle=\ttfamily\scriptsize, frame=none,
        backgroundcolor=\color{gray!8}, aboveskip=2pt, belowskip=2pt}

\decktitle{5}{A6}{Worked Example: The RBC Model by Perturbation}{Thursday, October 22, 2026}

\begin{document}

\makedecktitle

%=============================================================================
\begin{frame}{Where We Are}
\CourseMap{5}{5.A6}
\vspace{-1mm}
\begin{center}
\small \textbf{Block A:} 5.A1 what perturbation computes $\;\cdot\;$ 5.A2 the model and the parameter $\;\cdot\;$ 5.A3 first order $\;\cdot\;$ 5.A4 linear RE systems $\;\cdot\;$ 5.A5 higher order, risk, pruning $\;\cdot\;$ \textbf{5.A6} the worked example.
\end{center}
\end{frame}

%=============================================================================
\begin{frame}{The Same Model, a Third Way}
\begin{columns}[T]
\begin{column}{0.53\textwidth}
\small
3.A4 solved this model with every choice made the crudest way. 4.B5 solved it with the Session 4 toolkit. Both put values on a grid and iterated to a fixed point.
\medskip

This deck solves it a third way, and nothing about the procedure resembles the first two: no grid, no iteration, no maximization. Five steps, all of them linear algebra.
\medskip

\begin{center}\footnotesize
\begin{tabularx}{\linewidth}{@{}>{\raggedright\arraybackslash}p{1.6cm} >{\raggedright\arraybackslash}X@{}}
\toprule
\textbf{4.B5} & \textbf{here} \\
\midrule
grid in $k$ & \textbf{one point}, $k^{*}$ \\
Rouwenhorst $\Pi$ & \textbf{$\rho$ and $\sigma$ directly} \\
golden section & \textbf{no maximization at all} \\
Howard sweeps & \textbf{one QZ decomposition} \\
\bottomrule
\end{tabularx}
\end{center}
\end{column}
\begin{column}{0.44\textwidth}
\begin{conceptbox}[title=\textbf{The model}]\small
\[
  c+k'=zk^{\alpha}+(1-\delta)k,\quad \ln z'=\rho\ln z+\sigma\varepsilon'
\]
with $\ln$ utility. Quarterly: $\alpha=0.33$, $\beta=0.99$, $\delta=0.025$, $\rho=0.95$, $\sigma=0.007$.
\end{conceptbox}
\medskip
\begin{examplebox}\footnotesize
$\delta\neq1$, so there is no closed-form policy. The benchmarks are the steady state, the Euler residual, and 4.B5's own solution.
\end{examplebox}
\end{column}
\end{columns}
\end{frame}

%=============================================================================
\begin{frame}{What ``Perturbation'' Commits Us To}
\begin{columns}[T]
\begin{column}{0.53\textwidth}
\small
Six decisions, and they are different decisions from the six of 3.A2:
\medskip
\begin{center}\footnotesize
\begin{tabularx}{\linewidth}{@{}>{\raggedright\arraybackslash}p{2.5cm} >{\raggedright\arraybackslash}X@{}}
\toprule
\textbf{Choice} & \textbf{Here} \\
\midrule
expansion point & deterministic steady state \\
variable & log deviations \\
order & first, then second \\
root selection & stable, by Blanchard--Kahn \\
derivatives & numerical (Dynare would use analytic) \\
domain & \emph{unspecified} --- and that is the problem \\
\bottomrule
\end{tabularx}
\end{center}
\end{column}
\begin{column}{0.44\textwidth}
\begin{cautionbox}[title=\textbf{The last row}]\small
A grid method tells you its domain: it is the grid. Perturbation does not. Nothing in the output says where the answer stops being usable --- which is exactly what we will have to measure.
\end{cautionbox}
\medskip
\begin{takeawaybox}\footnotesize
So the deck's structure is forced: solve in three slides, then spend four finding out where the solution is valid.
\end{takeawaybox}
\end{column}
\end{columns}
\end{frame}

%=============================================================================
\begin{frame}[fragile]{Step 1: The Steady State, and the Linearized System}
\begin{columns}[T]
\begin{column}{0.56\textwidth}
\begin{lstlisting}
kss = ((1/beta - 1 + delta)/alpha)**(1/(alpha-1))
yss = kss**alpha                 # 3.0153
css = yss - delta*kss            # 2.3066
kappa = beta*alpha*kss**(alpha-1) # 0.03475

#  A E_t s' = B s,  s = (khat, z, chat)
A = np.array([[kss, 0.0, 0.0],
              [0.0, 1.0, 0.0],
              [-kappa*(alpha-1), -kappa, 1.0]])
B = np.array([[alpha*yss + (1-delta)*kss, yss, -css],
              [0.0, rho, 0.0],
              [0.0, 0.0, 1.0]])
\end{lstlisting}
\end{column}
\begin{column}{0.41\textwidth}
\begin{conceptbox}[title=\textbf{Check it first}]\small
$k^{*}=28.3484$ --- identical to 4.B5's, to every digit. Two completely different programs agreeing on the one quantity they share is the cheapest test available.
\end{conceptbox}
\medskip
\begin{examplebox}\footnotesize
Rows of $A$ and $B$: resource constraint, shock process, Euler equation. Every entry is an elasticity evaluated at the steady state.
\end{examplebox}
\end{column}
\end{columns}
\vspace{1mm}
\src{Ledger \texttt{pert\_first\_order}. $k^{*}=28.348419$; 4.B5 reports $28.3484$.}
\end{frame}

%=============================================================================
\begin{frame}[fragile]{Step 2: One QZ Decomposition}
\begin{columns}[T]
\begin{column}{0.56\textwidth}
\begin{lstlisting}
def stable(a, b):
    return np.abs(b) < np.abs(a)   # |omega| < 1

S, T, aa, bb, Q, Z = ordqz(A, B, sort=stable,
                           output="real")
omega = np.abs(bb)/np.abs(aa)
n_unstable = np.sum(omega > 1.0)   # 1

n_x = 2                            # khat, z
Z11, Z21 = Z[:n_x,:n_x], Z[n_x:,:n_x]
F = Z21 @ inv(Z11)                 # chat  = F x
P = Z11 @ solve(S[:n_x,:n_x],
                T[:n_x,:n_x]) @ inv(Z11)
\end{lstlisting}
\end{column}
\begin{column}{0.41\textwidth}
\begin{conceptbox}[title=\textbf{The whole solve}]\small
One library call, then two matrix products. $|\omega|=0.950,\,0.962,\,\mathbf{1.050}$: one root outside the unit circle, one jump variable. Blanchard--Kahn holds.
\end{conceptbox}
\medskip
\begin{cautionbox}\footnotesize
If the count were wrong, \texttt{inv(Z11)} would fail --- the algorithm cannot form a solution that does not exist. Do not wrap this in a \texttt{try}.
\end{cautionbox}
\end{column}
\end{columns}
\vspace{1mm}
\src{Ledger \texttt{blanchard\_kahn}. \texttt{scipy.linalg.ordqz}; Klein (2000).}
\end{frame}

%=============================================================================
\begin{frame}{Step 3: Read Off the Policy, and the Impulse Response}
\begin{columns}[T]
\begin{column}{0.45\textwidth}
\small
The solution is four numbers:
\medskip
\begin{center}\footnotesize
\begin{tabularx}{\linewidth}{@{}>{\raggedright\arraybackslash}X r@{}}
\toprule
$\hat{k}'$ on $\hat{k}$ & $0.962061$ \\
$\hat{k}'$ on $z$ & $0.080097$ \\
$\hat{c}$ on $\hat{k}$ & $0.590408$ \\
$\hat{c}$ on $z$ & $0.322850$ \\
\bottomrule
\end{tabularx}
\end{center}
\medskip
Everything else follows. The impulse response to a one-standard-deviation innovation is a loop over 40 quarters --- no simulation, no Monte Carlo.
\end{column}
\begin{column}{0.52\textwidth}
\begin{center}
\begin{tikzpicture}
\begin{axis}[
    width=7.4cm, height=5.0cm,
    xmin=0, xmax=39, ymin=-0.02, ymax=0.75,
    xlabel={\scriptsize quarters}, ylabel={\scriptsize \% deviation},
    tick label style={font=\scriptsize}, label style={font=\scriptsize},
    scaled ticks=false, /pgf/number format/fixed,
    axis x line*=bottom, axis y line*=left,
    legend style={font=\tiny, draw=none, fill=none,
                  at={(0.98,0.97)}, anchor=north east}]
  \addplot[very thick, black!45] coordinates {
      (0,0.7000) (1,0.6650) (2,0.6317) (3,0.6002) (4,0.5702) (5,0.5416) (6,0.5146) (7,0.4888) 
      (8,0.4644) (9,0.4412) (10,0.4191) (11,0.3982) (12,0.3783) (13,0.3593) (14,0.3414) (15,0.3243) 
      (16,0.3081) (17,0.2927) (18,0.2781) (19,0.2641) (20,0.2509) (21,0.2384) (22,0.2265) 
      (23,0.2151) (24,0.2044) (25,0.1942) (26,0.1845) (27,0.1752) (28,0.1665) (29,0.1582) 
      (30,0.1502) (31,0.1427) (32,0.1356) (33,0.1288) (34,0.1224) (35,0.1163) (36,0.1104) 
      (37,0.1049) (38,0.0997) (39,0.0947) };
  \addlegendentry{$z$}
  \addplot[very thick, RedTitle] coordinates {
      (0,0.0000) (1,0.0561) (2,0.1072) (3,0.1537) (4,0.1960) (5,0.2342) (6,0.2687) (7,0.2997) 
      (8,0.3275) (9,0.3523) (10,0.3743) (11,0.3936) (12,0.4106) (13,0.4253) (14,0.4380) (15,0.4487) 
      (16,0.4576) (17,0.4649) (18,0.4708) (19,0.4752) (20,0.4783) (21,0.4802) (22,0.4811) 
      (23,0.4810) (24,0.4800) (25,0.4782) (26,0.4756) (27,0.4723) (28,0.4684) (29,0.4640) 
      (30,0.4590) (31,0.4537) (32,0.4479) (33,0.4418) (34,0.4353) (35,0.4286) (36,0.4217) 
      (37,0.4145) (38,0.4072) (39,0.3997) };
  \addlegendentry{$\hat{k}$}
  \addplot[very thick, orange!85!black] coordinates {
      (0,0.2260) (1,0.2478) (2,0.2673) (3,0.2845) (4,0.2998) (5,0.3132) (6,0.3248) (7,0.3348) 
      (8,0.3433) (9,0.3504) (10,0.3563) (11,0.3609) (12,0.3645) (13,0.3671) (14,0.3688) (15,0.3696) 
      (16,0.3697) (17,0.3690) (18,0.3677) (19,0.3658) (20,0.3634) (21,0.3605) (22,0.3572) 
      (23,0.3535) (24,0.3494) (25,0.3450) (26,0.3403) (27,0.3354) (28,0.3303) (29,0.3250) 
      (30,0.3195) (31,0.3139) (32,0.3082) (33,0.3024) (34,0.2965) (35,0.2906) (36,0.2846) 
      (37,0.2786) (38,0.2726) (39,0.2666) };
  \addlegendentry{$\hat{c}$}
\end{axis}
\end{tikzpicture}
\end{center}
\end{column}
\end{columns}
\vspace{1mm}
\src{Ledger \texttt{fig\_irf} and \texttt{pert\_first\_order}.}
\end{frame}

%=============================================================================
\begin{frame}{Reading the Impulse Response}
\begin{columns}[T]
\begin{column}{0.55\textwidth}
\small
The shock itself decays at $\rho=0.95$: half of it is gone after $13.5$ quarters.
\medskip

Consumption jumps immediately, by $0.226\%$, and keeps rising to a peak of $\mathbf{0.370\%}$ --- it does not track the shock, because capital is still accumulating.
\medskip

Capital starts at zero, by construction: it is predetermined. It peaks at $\mathbf{0.481\%}$ in quarter $\mathbf{22}$ --- five and a half years after a shock that was half gone in three.
\medskip

That lag is $\kappa=0.03475$ doing its work, and it is the classic RBC propagation mechanism: the shock is transitory, the capital stock's response is not.
\end{column}
\begin{column}{0.42\textwidth}
\begin{takeawaybox}\footnotesize
None of this needed a simulation. From four numbers, the entire dynamic response --- and, with two more lines, every second moment in closed form.
\end{takeawaybox}
\medskip
\begin{cautionbox}\footnotesize
And every one of these numbers is \textbf{certainty equivalent}. Redo it with $\sigma$ ten times larger and the picture is identical.
\end{cautionbox}
\end{column}
\end{columns}
\vspace{1mm}
\src{Ledger \texttt{pert\_first\_order}: peak $\hat{c}$ $0.3697\%$, peak $\hat{k}$ $0.4811\%$ at quarter $22$, shock half-life $13.5$ quarters.}
\end{frame}

%=============================================================================
\begin{frame}{Step 4: Grade It, Across the Whole State Space}
\begin{columns}[T]
\begin{column}{0.52\textwidth}
\begin{center}
\begin{tikzpicture}
\begin{axis}[
    width=7.7cm, height=5.4cm,
    xmin=14, xmax=43, ymin=-4.6, ymax=-2.3,
    xlabel={\scriptsize capital $k$}, ylabel={\scriptsize $\log_{10}\max|\mathcal{E}|$},
    tick label style={font=\scriptsize}, label style={font=\scriptsize},
    scaled ticks=false, /pgf/number format/fixed,
    /pgf/number format/precision=1,
    axis x line*=bottom, axis y line*=left]
  \addplot[very thick, RedTitle] coordinates {
      (14.174,-2.554) (14.655,-2.595) (15.135,-2.636) (15.616,-2.678) (16.096,-2.720) 
      (16.577,-2.762) (17.057,-2.804) (17.538,-2.847) (18.018,-2.891) (18.499,-2.935) 
      (18.979,-2.980) (19.460,-3.026) (19.940,-3.073) (20.420,-3.121) (20.901,-3.170) 
      (21.381,-3.221) (21.862,-3.273) (22.342,-3.327) (22.823,-3.382) (23.303,-3.440) 
      (23.784,-3.501) (24.264,-3.565) (24.745,-3.632) (25.225,-3.703) (25.706,-3.780) 
      (26.186,-3.862) (26.667,-3.952) (27.147,-4.051) (27.628,-4.163) (28.108,-4.293) 
      (28.589,-4.315) (29.069,-4.193) (29.550,-4.092) (30.030,-4.006) (30.511,-3.930) 
      (30.991,-3.862) (31.472,-3.802) (31.952,-3.746) (32.433,-3.696) (32.913,-3.649) 
      (33.393,-3.606) (33.874,-3.566) (34.354,-3.528) (34.835,-3.493) (35.315,-3.460) 
      (35.796,-3.429) (36.276,-3.399) (36.757,-3.371) (37.237,-3.344) (37.718,-3.319) 
      (38.198,-3.295) (38.679,-3.272) (39.159,-3.250) (39.640,-3.228) (40.120,-3.208) 
      (40.601,-3.188) (41.081,-3.170) (41.562,-3.152) (42.042,-3.134) (42.523,-3.117) };
  \draw[black!45, dashed] (axis cs:14,-2.514) -- (axis cs:43,-2.514);
  \node[font=\tiny, text=black!60, anchor=south west] at (axis cs:15.0,-2.50)
    {global VFI, worst on the whole grid};
  \draw[black!35, dotted, thick] (axis cs:28.348,-4.6) -- (axis cs:28.348,-2.3);
  \node[font=\tiny, text=black!55, anchor=south, rotate=90] at (axis cs:28.0,-4.3) {$k^{*}$};
\end{axis}
\end{tikzpicture}
\end{center}
\end{column}
\begin{column}{0.45\textwidth}
\small
The Euler residual of the first-order solution, at every capital level in 4.B5's grid.
\medskip

It is a \textbf{V}, centred on $k^{*}$, and it is exactly the shape the theory predicts: best where we expanded, worse in both directions, and monotone in the distance.
\medskip

\begin{center}\footnotesize
\begin{tabularx}{\linewidth}{@{}>{\raggedright\arraybackslash}X r@{}}
\toprule
at $k^{*}$ & $10^{-4.32}$ \\
at $k=0.5\,k^{*}$ & $10^{-2.55}$ \\
\midrule
accuracy lost & \textbf{1.8 decades} \\
\bottomrule
\end{tabularx}
\end{center}
\end{column}
\end{columns}
\vspace{1mm}
\src{Ledger \texttt{fig\_pert\_error} and \texttt{three\_methods}. ``One decade'' is a factor of ten in the error (3.A1).}
\end{frame}

%=============================================================================
\begin{frame}{The Uncomfortable Comparison}
\begin{columns}[T]
\begin{column}{0.5\textwidth}
\small
Now put the two solutions side by side on the same model, the same grid, the same residual definition.
\medskip
\begin{center}\footnotesize
\begin{tabularx}{\linewidth}{@{}>{\raggedright\arraybackslash}X r r@{}}
\toprule
 & \textbf{4.B5 grid} & \textbf{1st order} \\
\midrule
numbers stored & $2{,}100$ & $\mathbf{4}$ \\
solve time & $0.94$\,s & $<0.001$\,s \\
worst $\mathcal{E}$, near $k^{*}$ & $10^{-2.82}$ & $\mathbf{10^{-4.01}}$ \\
worst $\mathcal{E}$, whole grid & $10^{-2.51}$ & $10^{-2.56}$ \\
\bottomrule
\end{tabularx}
\end{center}
\medskip
Perturbation is $1.2$ decades \textbf{better} near the steady state, and --- read it carefully --- it is \textbf{not worse anywhere}. Its worst residual on the whole grid, $10^{-2.56}$, essentially matches the grid solver's $10^{-2.51}$.
\end{column}
\begin{column}{0.47\textwidth}
\begin{cautionbox}[title=\textbf{The wrong conclusion}]\small
``So perturbation is at least as good as a global method.'' No. What this shows is that \textbf{the grid solver is not a gold standard}. Its own error is $10^{-2.5}$, dominated by interpolation and the grid --- so it cannot certify anything finer than that.
\end{cautionbox}
\medskip
\begin{takeawaybox}\footnotesize
Comparing two approximations tells you they differ. It does not tell you which is right. You need a \emph{third}, better one --- and Block B builds it.
\end{takeawaybox}
\end{column}
\end{columns}
\vspace{1mm}
\src{Ledger \texttt{three\_methods}. Both rows use the same unit-free residual on the same random draws.}
\end{frame}

%=============================================================================
\begin{frame}{What One More Block of Algebra Buys}
\begin{columns}[T]
\begin{column}{0.5\textwidth}
\small
Solve the twelve second-order conditions of 5.A5 and re-grade:
\medskip
\begin{center}\footnotesize
\begin{tabularx}{\linewidth}{@{}>{\raggedright\arraybackslash}X r r@{}}
\toprule
 & \textbf{1st} & \textbf{2nd} \\
\midrule
numbers stored & $4$ & $12$ \\
solve time & $<0.001$\,s & $0.02$\,s \\
worst $\mathcal{E}$ near $k^{*}$ & $10^{-4.01}$ & $\mathbf{10^{-6.14}}$ \\
worst $\mathcal{E}$, whole grid & $10^{-2.56}$ & $\mathbf{10^{-4.50}}$ \\
\bottomrule
\end{tabularx}
\end{center}
\medskip
\textbf{Roughly two decades, on both bands, for eight more numbers.} That is the best return on effort anywhere in Sessions 3 to 5 --- and it also buys the risk correction, which first order reports as identically zero.
\end{column}
\begin{column}{0.47\textwidth}
\begin{takeawaybox}\footnotesize
Second-order perturbation now beats the $2{,}100$-number grid solution \textbf{everywhere on the grid}, by two decades, using twelve numbers.
\end{takeawaybox}
\medskip
\begin{examplebox}\footnotesize
If the state is one-dimensional, smooth, and you care about behaviour near the steady state, this is very hard to beat. All three conditions matter.
\end{examplebox}
\end{column}
\end{columns}
\vspace{1mm}
\src{Ledger \texttt{three\_methods}. Second-order coefficients from the 12-equation residual-matching solve of 5.A5, residual $3.6\times10^{-9}$.}
\end{frame}

%=============================================================================
\begin{frame}{What We Deliberately Did Not Do}
\begin{columns}[T]
\begin{column}{0.53\textwidth}
\small
\begin{itemize}\small
  \item \textbf{No occasionally binding constraint.} Put a borrowing limit in and the policy has a kink; the expansion does not exist at it (5.A1, 5.A5).
  \item \textbf{No large shocks.} $\sigma=0.007$ per quarter. 5.A5 measured the $\sigma^{2}$ law failing by $60\%$ at twenty times that.
  \item \textbf{No second endogenous state.} Perturbation scales to fifty states far better than a grid does --- its real advantage, and this model cannot show it.
  \item \textbf{No global check.} We graded against residuals and against 4.B5, neither of which is the truth.
  \item \textbf{No analytic derivatives.} Dynare uses them; ours are numerical, which cost six orders of magnitude of residual at second order until the step size was fixed.
\end{itemize}
\end{column}
\begin{column}{0.44\textwidth}
\begin{cautionbox}[title=\textbf{The honest summary}]\small
Perturbation is extraordinarily good at what it does and gives you \textbf{no warning at all} when you leave the region where it does it.
\end{cautionbox}
\medskip
\begin{examplebox}\footnotesize
Which is the argument for Block B: not that projection is more accurate here --- it is --- but that it tells you its own domain.
\end{examplebox}
\end{column}
\end{columns}
\end{frame}

%=============================================================================
\begin{frame}{Takeaways}
\begin{takeawaybox}
\begin{itemize}\small
  \item The same model as 3.A4 and 4.B5, solved with \textbf{no grid, no iteration and no maximization} --- a steady state, two matrices, and one QZ decomposition.
  \item The entire first-order solution is \textbf{four numbers}, and $k^{*}=28.3484$ agrees with 4.B5 to every digit.
  \item From those four numbers: the full impulse response. Capital peaks at $0.481\%$ in quarter \textbf{22}, from a shock that is half gone by quarter $13.5$.
  \item The residual is a \textbf{V centred on $k^{*}$}: $10^{-4.32}$ there, $10^{-2.55}$ at half the steady-state capital. \textbf{1.8 decades} lost, exactly as the theory says.
  \item Against 4.B5 it is \textbf{1.2 decades better near $k^{*}$ and no worse anywhere} --- which says less about perturbation than about the grid solver, whose own error is $10^{-2.5}$.

  \item \textbf{Second order buys two more decades for eight more numbers}, and beats the $2{,}100$-number grid solution everywhere on it. Comparing two approximations cannot tell you which is right --- Block B builds the third one.
\end{itemize}
\end{takeawaybox}
\end{frame}

%=============================================================================
\begin{frame}{Bridge: Keep the Accuracy, Lose the Grid}
\begin{columns}[T]
\begin{column}{0.53\textwidth}
\small
Block A traded \textbf{global validity} for speed and dimension, and we measured exactly what the trade cost: nothing near the steady state, $1.8$ decades at half of it, and everything at a kink.
\medskip

Block B makes the opposite trade. Keep the solution \textbf{global} --- valid on a whole domain, constraints and all --- but stop storing one number per grid point. Store a handful of coefficients instead, and define ``solved'' as ``the residual is zero''.
\end{column}
\begin{column}{0.44\textwidth}
\begin{conceptbox}[title=\textbf{After the break (5.B1)}]\small
Projection: the residual function, and what it means to make a function small. Then bases, Chebyshev, finite elements, collocation against Galerkin --- and the same model again, solved a fourth way.
\end{conceptbox}
\medskip
\begin{examplebox}\footnotesize
Spoiler, from the ledger: $56$ numbers, and \textbf{four decades} better than anything in this block.
\end{examplebox}
\end{column}
\end{columns}
\end{frame}

\end{document}
